\chapter[Appendix: Source Code]{Source Code and Reproducibility}
\section{Source Snapshot}
The \texttt{source/} directory contains unmodified copies of the project's five Python files. The \texttt{evidence/manifest.json} file records SHA-256 hashes of the source and workbook, allowing the code used for verification to be checked. Comments starting with \texttt{ponytail:} mark deliberate simplifications and name their known limit. The following excerpts are extracted automatically from that snapshot rather than rewritten to improve the reported outcomes.
\par\needspace{10\baselineskip}\begingroup\footnotesize\setlength{\tabcolsep}{4pt}
\begin{longtable}{>{\raggedright\arraybackslash}p{\dimexpr 0.22000\linewidth-2\tabcolsep\relax}>{\raggedright\arraybackslash}p{\dimexpr 0.15000\linewidth-2\tabcolsep\relax}>{\raggedright\arraybackslash}p{\dimexpr 0.63000\linewidth-2\tabcolsep\relax}}
\caption{Bundled source files}\\
\toprule
\textbf{File} & \textbf{Lines} & \textbf{Snapshot path} \\ \midrule
\endfirsthead
\multicolumn{3}{l}{\small\itshape Bundled source files (continued)}\\
\toprule
\textbf{File} & \textbf{Lines} & \textbf{Snapshot path} \\ \midrule
\endhead
\midrule
\multicolumn{3}{r}{\scriptsize Continued on the next page}\\
\endfoot
\bottomrule\endlastfoot
checks.py & 1284 & source/checks.py \\
live.py & 192 & source/live.py \\
parking.py & 309 & source/parking.py \\
rm.py & 339 & source/rm.py \\
sim.py & 651 & source/sim.py \\
\end{longtable}\endgroup

\section{Core Code Excerpts}
\subsection*{Legal state transitions}
\noindent\textit{Source: sim.py}
\VerbatimInput[fontsize=\scriptsize,baselinestretch=1,frame=single,numbers=left,numbersep=4pt]{generated/snippet-1.txt}
\subsection*{Authentication and account lockout}
\noindent\textit{Source: sim.py}
\VerbatimInput[fontsize=\scriptsize,baselinestretch=1,frame=single,numbers=left,numbersep=4pt]{generated/snippet-2.txt}
\subsection*{Identifier-based test runner}
\noindent\textit{Source: checks.py}
\VerbatimInput[fontsize=\scriptsize,baselinestretch=1,frame=single,numbers=left,numbersep=4pt]{generated/snippet-3.txt}
\subsection*{95th-percentile measurement}
\noindent\textit{Source: checks.py}
\VerbatimInput[fontsize=\scriptsize,baselinestretch=1,frame=single,numbers=left,numbersep=4pt]{generated/snippet-4.txt}

\clearpage
\section{Compilation and Reproduction}
Compile \texttt{main.tex} three times with XeLaTeX (or run \texttt{build.ps1}) to update the contents, figure list, and references. Figures are drawn directly in LaTeX with TikZ, so rendering them needs no external service or network connection. Font fallbacks support Overleaf.

Reproducing the evidence requires Python and \texttt{openpyxl}. Run the following from the report directory:
\par\noindent\begin{minipage}{\linewidth}
\begin{Verbatim}[fontsize=\small,frame=single]
python refresh_evidence.py
xelatex -interaction=nonstopmode -halt-on-error main.tex
xelatex -interaction=nonstopmode -halt-on-error main.tex
xelatex -interaction=nonstopmode -halt-on-error main.tex
\end{Verbatim}
\end{minipage}\par\medskip\noindent The first command refreshes generated tables and measurements; timing values may vary by machine. It does not write to the original project workbook. Detailed instructions, directory layout, and editable cover fields are described in the accompanying \texttt{README.md}.
